\begin{abstract}
Instant messaging (IM) platforms have become primary channels for distributing time-sensitive information in both academic and organizational contexts. Despite this, productivity tooling has not kept pace: users are left to manually identify, extract, and transfer actionable tasks from raw chat streams into structured scheduling applications, a process that is both error-prone and cognitively taxing. We present \textbf{Chatnalyxer}, a fully passive, multi-modal pipeline that processes WhatsApp messages at the network edge, extracts temporal tasks using a Large Language Model (LLM), and schedules them automatically without any user interaction. A key engineering decision—decoupling the Node.js edge ingestion layer from the FastAPI processing core through an asynchronous HTTP 202 WebHook—bounds local ingestion latency to $\mathcal{O}(1)$ under burst traffic, mitigating the TCP socket exhaustion that plagues synchronous architectures. To handle the deictic temporal language prevalent in chat (e.g., ``tomorrow at 5''), we embed a real-time UTC timestamp into each LLM prompt, enabling precise ISO-8601 deadline resolution without any model fine-tuning. Non-text media—voice notes and images—are normalized into UTF-8 through Azure Cognitive Services (ASR and OCR, respectively) before reaching the inference layer. Evaluated against the Chatnalyxer Evaluation Corpus (CEC-10K), the pipeline achieves an F1 score of 0.94 on plain-text extraction, demonstrating that ambient, in-stream productivity management is architecturally feasible.
\end{abstract}
